\subsection{The relative lift and the identity of the witness}

\begin{result}{A.2}{Proved}{EXIT\_TOWER E.46 is confirmed on the dependencies of E.49}
The relative lifting statement in EXIT\_TOWER E.46 has the same actual source and prescribed-conjugator conclusion as the refereed relative lifting construction used in LM7\_PUSH E.58. Choosing the same relative centralizer lift and the same compact lifting point makes their witnesses identical. Without those choices, each text specifies a nonempty family, rather than one canonical automorphism. Every member of the stated family has
\[
 \lambda=1,\quad \rho=0,\quad
 T=5I\pmod {12},\quad W=I\pmod {12},\quad\Delta=1.
 \tag{A.L1}
\]
No step of EXIT\_TOWER E.47 is used in the proof of E.49. Its procongruence multitwist-normalizer assertion is therefore not an input to the main decision here.
\end{result}

\paragraph{Typed lifting data.}
Let $\mathcal B=\widehat{C_2(s)*C_3(r)}$, $\tau=sr$, and
$F=\ker(\mathcal B\twoheadrightarrow\mathfrak S_3)$, with
$x=\tau^{-2}$ and $y=sxs$. All subgroups in this paragraph are closed.
Put $H=F'F^3$ and let
\[
 R_H=\bigcap_{H\twoheadrightarrow L,\ L\text{ finite nilpotent}}
              \ker(H\twoheadrightarrow L).
\]
Both $F$ and $H$ are characteristic in $\mathcal B$. The first assertion follows because the epimorphisms to $\mathfrak S_3$ form one orbit under $\Aut(\mathfrak S_3)$ and hence have one kernel. The second then follows from its definition. The subgroup $R_H$ is characteristic in $H$ and normal in $\mathcal B$.

The natural quotient maps give isomorphisms of profinite groups
\[
 H/R_H\simeq\prod_p H^{(p)},\qquad
 \mathcal B/R_H\simeq
       \prod_{p,\,\mathcal B/H}\mathcal B/O^p(H).
 \tag{A.L2}
\]
For the first map, every finite collection of the distinct-prime factors is jointly onto: a common quotient of finite groups of coprime orders is trivial, and induction using the subgroup version of Goursat's lemma gives surjectivity onto the product. Compactness gives surjectivity onto the infinite product, and the kernel is the displayed nilpotent residual. For the second map, lift the common $\mathcal B/H$ coordinate by one element of $\mathcal B$ and lift the remaining discrepancy by the first isomorphism. These constructions use no splitting of $\mathcal B\to\mathcal B/H$.

\paragraph{Audit of the finite lifting calculation.}
We give the points at which E.46 needs more than an ordinary statement about conjugacy classes. Take a finite quotient $\Lambda\twoheadrightarrow E$ of the abstract source $\Lambda=C_2*C_3$, a normal subgroup $A\triangleleft E$, and $D=E/A$. The classes prescribed at the three peripherals are $A$-classes. For a subgroup $J\leq E$ mapping onto $D$, put $M=J\cap A$. For every $J$-invariant irreducible complex representation $\rho$ of $M$, of dimension $d$, let $e=\exp J$ and form
\[
 \widetilde J_\rho=
 \{(g,T):T\rho(m)T^{-1}=\rho(gmg^{-1}),\ (\det T)^e=1\}.
 \tag{A.L3}
\]
Schur's lemma leaves exactly $de$ implementers over each $g$. The subgroup
$m\mapsto(m,\rho(m))$ is normal and intersects the scalar kernel trivially. Its quotient $Q_\rho$ is a finite central extension of $D$, and the map
\[
 \widetilde J_\rho\longrightarrow J\times_D Q_\rho,
       \qquad (g,T)\longmapsto(g,[(g,T)])
\]
is injective and has equal finite source and target orders, hence is an isomorphism. This proves the precise pullback used by the trace computation.

The order-two and order-three source generators lift to $Q_\rho$ with their actual orders. An implementer's square or cube is scalar, and rescaling it over $\mathbb C$ makes the corresponding power one. Its determinant then satisfies the condition in (A.L3). Taking the joint image of these source maps gives the central extension $D^\dagger\to D$ used in E.46. It is an image of the actual source; no abstract central extension is assumed to split.

For the two prescribed systems of classes, write $C_i^j\subset J$, $j=0,1$, and let $L_i^j$ be their implementer sums. Each sum commutes with $\rho(M)$ and is scalar. Only the $J$-invariant irreducibles contribute: a class sum supported on a coset moving an irreducible idempotent has zero product with that idempotent. The exact homomorphism count is
\[
 H_j(J)=\frac1{|M|}
   \sum_{[\rho]\text{ invariant under }J}
       \frac{\operatorname{tr}(L_1^j)
             \operatorname{tr}(L_2^j)
             \operatorname{tr}(L_3^j)}{\dim\rho}.
 \tag{A.L4}
\]
Indeed the regular-character term is
$d\operatorname{tr}(L_1^jL_2^jL_3^j)$, and each scalar sum has scalar value $\operatorname{tr}(L_i^j)/d$. This gives the denominator $d$ in (A.L4).

Choose in $J$ lifts of the three prescribed $D$ coordinates. Their difference from the chosen lifts in $E$ lies in $A$, so conjugation identifies the two class sets. In the implementer pullback the corresponding matrices conjugate $L_i^0$ to $L_i^1$; all three traces agree. Thus $H_0(J)=H_1(J)$ for every $J$. Subgroups not mapping onto $D$ contribute zero to both counts. M\"obius inversion on the finite subgroup poset gives equality of epimorphism counts. The original quotient map is counted on one side, so an epimorphism with the required three conjugator coordinates exists on the other.

The relative compactness step is also exact. For a finite family of quotient and generating conditions, take one finite source quotient $E_0$ containing them and the quotient by $H$. Replace the image of $H$ in it by its nilpotent quotient, obtaining $D_0$. The central extension $D^\dagger$ constructed above still has nilpotent image of $H$, since a central extension of a nilpotent finite group is nilpotent. Its source map therefore factors through $\mathcal B/R_H$. Precomposition with the prescribed relative automorphism supplies the second system of peripheral data. The positive count provides a point meeting this finite family of conditions. The compact space of three conjugators has the finite-intersection property, giving one global point. Surjectivity on every finite quotient makes the resulting endomorphism onto. A surjective endomorphism of a finitely generated profinite group is injective: pullback permutes the finite set of open subgroups of each fixed index, and its kernel is consequently in every open subgroup. This checks the passage from a relative automorphism to an actual automorphism, with all three coordinates retained.

\paragraph{The retained finite frames.}
Let $q_{12}:\mathcal B\twoheadrightarrow
Q=\operatorname{SL}_2(\mathbb Z/12)/\{\pm I\}$ have the signed representatives
\[
 s\longmapsto S=\begin{pmatrix}0&1\\-1&0\end{pmatrix},
 \qquad
 r\longmapsto R_0=\begin{pmatrix}0&-1\\1&-1\end{pmatrix}.
\]
The free subgroup $F$ has rank two and $[\mathcal B:H]=54$.
Its signed image is $V_4\times\operatorname{SL}_2(\mathbb F_3)$ of order $96$; commutators and cubes give
\[
 q_{12}(H)=V_4\times Q_8,\qquad
 q_{12}([x,y]^2)=[5I].
 \tag{A.L5}
\]
The joint image in $Q\times\mathcal B/H$ is the fibre product over $Q/q_{12}(H)$ and has order $1728$. Surjectivity onto that fibre product follows by correcting a lift of a prescribed $\mathcal B/H$ coordinate using $H$. The word $[x,y]^2$ supplies its central element $([5I],1)$.

In the relative pro-two quotient $\mathcal B/O^2(H)$, coprime complement conjugacy lifts this centralizer element to an element $c$ centralizing $r$, lying in $H^{(2)}$, with the stated finite image. For clarity, the kernel of the map to the joint finite image is pro-two. The two order-three complements generated by $r$ and by a conjugate lifting a centralizer element are conjugate by that kernel. At finite level this follows by induction through elementary abelian central layers and the identity $\ker(1+a+a^2)=\operatorname{im}(a-1)$ over $\mathbb F_2$; compactness gives the profinite assertion.

Use the identity relative automorphism in (A.L2), with conjugator $c$ for $r$ in the two-adic component, identity in every other component, and identity conjugators for $s$ and $\tau^{-1}$. The lift just audited, normalized at $\tau$, has
\[
 \begin{gathered}
 \Phi(\tau)=\tau,\quad
 \Phi(s)=\beta s\beta^{-1},\quad
 \Phi(r)=ArA^{-1},\\
 \beta\in R_H,\quad A\in H,\quad
 q_{12}(\beta)=1,\quad q_{12}(A)=[5I].
 \end{gathered}
 \tag{A.L6}
\]
The element $f=s\Phi(s)$ lies in $R_H\subset F'$. The equations $\Phi(s)^2=1$ and $\Phi(r)^3=1$, using the conjugation actions of $s$ and $r$ on $(x,y,(xy)^{-1})$, give the two-cycle and three-cycle relations. The restriction to the characteristic subgroup $F$ is an automorphism. Thus this is an element of the full group $\GTz$.

The pure frames are unique, because the centralizers of $s$ and $r$ in $\mathcal B$ are their finite factors and these factors meet $F$ trivially. The profinite assertion follows directly from the finite-$p$ theorem applied to this virtually free group, or from its profinite free-product tree. Since $\beta\in F'$, its frame coordinate $\rho$ is zero. The signed lifts selected by the normalization at modulus four give exactly (A.L1): $5I$ is $I$ modulo four and $-I$ modulo three. These finite signed identities are checked in \texttt{check\_cm\_comparisons.py} and the descent certificate.

\paragraph{Relation to the previously named element.}
The notation $\Phi_0$ in LM7\_PUSH E.58 denotes a choice of the same $c$ and of a compact lifting point with (A.L6). Therefore that choice can be used verbatim in E.46. Conversely every choice allowed by E.46 satisfies the same finite frame equations and has the same character value. Two choices with the same relative data differ by an element acting trivially on that relative quotient and having zero difference of the normalized finite frame data. This does not identify the two actual automorphisms. In particular, existential notation alone supplies no unique canonical lift. The main proof below excludes every member of this family from $\GT$.
